\chapter{Conclusiones y Trabajo Futuro}

\section{Conclusiones}

El desarrollo de este Trabajo de Fin de Grado ha permitido diseñar e implementar un sistema de control para una pierna robótica bípeda de 3 grados de libertad utilizando controladores ODrive.

Los principales logros obtenidos son:

\begin{enumerate}
    \item \textbf{Estructura de Software Completa}: Se ha desarrollado un paquete de ROS2 (Jazzy) con todos los módulos necesarios para el control de la pierna robótica, incluyendo interfaces con ODrive, cinemática, y generación de trayectorias.
    
    \item \textbf{Integración con ODrive}: Se ha logrado la comunicación exitosa vía USB con el ODrive v3.6 oficial, implementando calibración automática y los tres modos de control (par, velocidad, posición).
    
    \item \textbf{Cinemática Validada}: Las implementaciones de cinemática directa e inversa han sido verificadas mediante tests unitarios, asegurando que la pierna pueda alcanzar posiciones objetivo en el espacio cartesiano.
    
    \item \textbf{Generación de Trayectorias}: Se han implementado algoritmos para generar trayectorias suaves de marcha, utilizando curvas de Bézier para evitar movimientos bruscos.
    
    \item \textbf{Cálculo de Torque}: Se han realizado los cálculos de par requerido para cada articulación en configuración de sentadilla, identificando que la articulación Hip 1 excede la capacidad del ODrive con reductora 16:1 (22.46 Nm requeridos vs 14.7 Nm disponibles).
\end{enumerate}

\section{Limitaciones Identificadas}

\begin{itemize}
    \item \textbf{ODrive Makerbase Bricked}: El clon chino Makerbase M1 M22015 no ha podido ser recuperado, limitando el sistema a 1 motor operativo en lugar de 3.
    
    \item \textbf{Capacidad de Par Insuficiente}: La articulación Hip 1 requiere más par del que puede proporcionar el ODrive con la reductora actual de 16:1, especialmente en configuraciones de sentadilla unipodal.
    
    \item \textbf{Encoders Pendientes}: La lectura de encoders AS5600 vía multiplexor I2C no ha sido probada físicamente debido a la falta de integración completa del hardware.
\end{itemize}

\section{Trabajo Futuro}

Las líneas de trabajo futuras para completar y mejorar el proyecto incluyen:

\subsection{Corto Plazo}
\begin{itemize}
    \item Recuperar o reemplazar el ODrive Makerbase para tener los 3 motores operativos por pierna.
    \item Probar la lectura de encoders AS5600 usando el multiplexor TCA9548A.
    \item Integrar la realimentación de posición de los encoders con el control de motores.
\end{itemize}

\subsection{Mediano Plazo}
\begin{itemize}
    \item Aumentar la relación de reducción de la articulación Hip 1 a 25:1 para satisfacer los requerimientos de par.
    \item Verificar la cinemática con la pierna física ensamblada (impresión 3D de los eslabones).
    \item Implementar control en lazo cerrado con realimentación de encoders.
\end{itemize}

\subsection{Largo Plazo}
\begin{itemize}
    \item Expandir el sistema a una configuración bípeda completa (2 piernas, 6 motores).
    \item Implementar algoritmos de marcha dinámica para caminata estable.
    \item Desarrollar interfaz de usuario para monitoreo y control en tiempo real.
    \item Integrar sensores de fuerza en la planta del pie para mejorar el control de equilibrio.
\end{itemize}

\section{Reflexión Final}

Este proyecto ha proporcionado una base sólida para el desarrollo de robots bípedos usando componentes de bajo costo como los controladores ODrive y motores D5065. A pesar de los desafíos con el hardware (Makerbase bricked), el sistema de software desarrollado es modular, escalable y listo para ser extendido a una plataforma robótica completa.
